\SourceLeaf{131}
Where seven or ten decimal places suffice, the ordinary tables or \textsc{Vlacq}'s, also published by \textsc{Vega}, are admittedly far more convenient than \textsc{Leonelli}'s. Indeed, without exception, the former tables are more accurate than necessary even in the most exacting astronomical calculations or other calculations relating to the physical world. In analytical calculations, however, or in the initial construction of tables, greater accuracy is often wanted. Then, provided one does not require more than 13 or 18 digits, \textsc{Leonelli}'s tables are unquestionably the most useful and convenient available for this purpose.

According to \textsc{Delambre}'s report, \textsc{Brigg} had in fact presented exactly the same method in his \emph{arithmetica logarithmica} and calculated a similar auxiliary table to 14 decimal places. Yet the French geometers themselves could find only one copy of this rare work, and that very part was missing from it. \textsc{Leonelli} thus retains the full credit for his independent invention and calculation, and it is only fair that we give him due thanks especially for the latter.

We have already stated the purpose of the second paper. We will now see how \textsc{Leonelli} proposes to achieve it. He proposes a table consisting of three columns. For brevity, instead of using \textsc{Leonelli}'s cumbersome and unnecessary terminology, we will denote corresponding entries in these three columns by $P$, $Q$, $R$. These quantities are to be related as follows. Let the number whose logarithm is $P$ be denoted by $x$. Then
\[
Q = \log\left(1+\frac{1}{x}\right), \qquad R = \log(1+x)
\]
Hence $R = P+Q$ always holds. The table is to be ordered using $x$ only through its logarithm $P$, so $P$ is to increase uniformly from zero to infinity, or rather until $Q$ may be regarded as vanishing. The use of such a table is easy to see. To determine from $\log a$ and $\log b$ the logarithm of $a+b$, assuming that $a$ is greater than $b$, enter the first column with $\log a - \log b$, or set this difference $= P$. Then clearly
\[
\log(a+b) = Q + \log a = R + \log b.
\]

